% ID: SEL-D0068-OA-E03
% ESTADO: cotejada. Segunda pasada visual y borrador compilado revisados.
% Fuente: Selectividad/sel_2009_matematicas_aplicadas/mat apli CC SS02.pdf
% SHA-256: 24ee6d89c206838cb0cf4f32c4e54391e7ca81f77c4dd0be651ad7218f37beb6
% Página: 1; opción A; ejercicio original 3.

\noindent\underline{Parte I}

Una enfermedad afecta al $10\%$ de la población. Una prueba de diagnóstico tiene las siguientes características: si se aplica a una persona con la enfermedad, da positivo en el $98\%$ de los casos; si se aplica a una persona que no tiene la enfermedad, da positivo en el $6\%$ de los casos. Se elige una persona, al azar, y se le aplica la prueba.
\begin{enumerate}[label=\alph*)]
\item \textbf{(1 punto)} \textquestiondown{}Cuál es la probabilidad de que dé positivo?
\item \textbf{(1 punto)} Si no da positivo, \textquestiondown{}cuál es la probabilidad de que tenga la enfermedad?
\end{enumerate}
\begin{samepage}
\noindent\underline{Parte II}

Se desea estimar la proporción de fumadores de una población mediante una muestra aleatoria.
\begin{enumerate}[label=\alph*)]
\item \textbf{(1 punto)} Si la proporción de fumadores en la muestra es $0.2$ y el error cometido en la estimación ha sido inferior a $0.03$, con un nivel de confianza del $95\%$, calcule el tamaño mínimo de la muestra.
\item \textbf{(1 punto)} Si en otra muestra de tamaño 280 el porcentaje de fumadores es del $25\%$, determine, para un nivel de confianza del $99\%$, el correspondiente intervalo de confianza para la proporción de fumadores de esa población.
\end{enumerate}
\end{samepage}
